% 原章第13—21页：6.2.1.3，式6.48—6.84及结论6.2。
\subsubsection{自编码器的参数估计}
\bookdeepheading{paragraph}{6.2.1.3.1}{学习带非负约束的自编码器}

现在介绍用于估计深度自编码器参数的反向传播算法（backpropagation algorithm）。施加非负约束，可以学习影像数据的稀疏、基于局部成分的表示。为在权重矩阵 $W$ 中引入非负约束，将函数 $C_W$ 改为以下二次约束（Hosseini-Asl 2016）：
\begin{equation}
C_W=\frac12\sum_{l=1}^2\sum_{i=1}^{s_l}\sum_{j=1}^{s_{l+1}}f(w_{ij}^{(l)}),
\tag{6.48}
\end{equation}
其中
\[
f(w_{ij})=\begin{cases}w_{ij}^2,&w_{ij}<0,\\0,&w_{ij}\geq0.\end{cases}
\]
学习带非负约束的自编码器，可表述为求解以下优化问题：
\begin{equation}
\min_{W,b}C_{\mathrm{SAE}}(W,b)=C_E(W,b)+\mu C_{\mathrm{KL}}(p\Vert\widehat p)+\lambda C_W.
\tag{6.49}
\end{equation}
反向传播算法（Ng 2011）可用于求解极小化问题 (6.49)。该算法源于随机梯度方法。令 $\theta=[W,b]$ 表示自编码器中的参数，$C_{\mathrm{SAE}}(\theta)=C_{\mathrm{SAE}}(W,b)$。更新目标函数 $C_{\mathrm{SAE}}(\theta)$ 的参数 $\theta$ 时，标准梯度下降算法为
\begin{equation}
\theta=\theta-\alpha\frac{\partial E[C_{\mathrm{SAE}}(\theta)]}{\partial\theta},
\tag{6.50}
\end{equation}
其抽样形式为
\begin{equation}
\theta=\theta-\alpha\frac{\partial C_{\mathrm{SAE}}(\theta)}{\partial\theta}.
\tag{6.51}
\end{equation}
具体而言，权重与偏置的更新规则分别为
\begin{equation}
w_{jk}^{(l)}=w_{jk}^{(l)}-\alpha\frac{\partial C_{\mathrm{SAE}}(W,b)}{\partial w_{jk}^{(l)}},
\tag{6.52}
\end{equation}
以及
\begin{equation}
b_j^{(l)}=b_j^{(l)}-\alpha\frac{\partial C_{\mathrm{SAE}}(W,b)}{\partial b_j^{(l)}},
\tag{6.53}
\end{equation}
其中，$\alpha>0$ 为学习率。目标函数关于权重的导数可进一步分解为
\begin{equation}
\frac{\partial C_{\mathrm{SAE}}(W,b)}{\partial w_{jk}^{(l)}}=
\frac{\partial C_E(W,b)}{\partial w_{jk}^{(l)}}+\mu\frac{\partial C_{\mathrm{KL}}(p\Vert\widehat p)}{\partial w_{jk}^{(l)}}+\lambda\gamma(w_{jk}^{(l)}),
\tag{6.54}
\end{equation}
其中
\[
\gamma(x)=\begin{cases}w_{jk},&w_{jk}<0,\\0,&w_{jk}\geq0.\end{cases}
\]
为方便起见，将式 (6.35) 改写为
\begin{equation}
C_E(W,b)=\frac1{2m}\sum_{u=1}^m C_E(W,b,x^{(u)}),
\tag{6.55}
\end{equation}
其中
\begin{equation}
C_E(W,b,x^{(u)})=\lVert y^{(u)}(W,b)-x^{(u)}\rVert^2,
\tag{6.56}
\end{equation}
以及
\begin{equation}
y^{(u)}(W,b)=a_j^{(2)}(W,b,x^{(u)})=\sigma\bigl(z_j^{(2)}(W,b,x^{(u)}).
\tag{6.57}
\end{equation}
由式 (6.54)、(6.55)，得到
\begin{equation}
\frac{\partial C_E(W,b)}{\partial w_{jk}^{(l)}}=\frac1{2m}\sum_{u=1}^m\frac{\partial C_E(W,b,x^{(u)})}{\partial w_{jk}^{(l)}}.
\tag{6.58}
\end{equation}
利用式 (6.30)、(6.31)、(6.56)、(6.57) 及链式法则，得到
\begin{equation}
\frac{\partial C_E(W,b,x^{(u)})}{\partial w_{jk}^{(2)}}=
-\bigl(y_j^u(W,b)-x_j^{(u)}\bigr)\sigma'(z_j^{(2)})a_k^{(1)}(x^{(u)}),
\tag{6.59}
\end{equation}
其中
\begin{equation}
\sigma'(x)=\sigma(x)(1-\sigma(x)).
\tag{6.60}
\end{equation}
令
\begin{equation}
\delta_j^{(2)}(x^{(u)})=\frac{\partial C_E(W,b,x^{(u)})}{\partial z_j^{(2)}}=
-(y_j^u(W,b)-x_j^{(u)})\sigma'(z_j^{(2)}).
\tag{6.61}
\end{equation}
将式 (6.61) 代入式 (6.59)，有
\begin{equation}
\frac{\partial C_E(W,b,x^{(u)})}{\partial w_{jk}^{(2)}}=\delta_j^{(2)}(x^{(u)})a_k^{(1)}(x^{(u)}),\quad j=1,\ldots,n,\ k=1,\ldots,n_1.
\tag{6.62}
\end{equation}
计算代价函数梯度的下一步，是沿网络将误差从最后一层反向传到隐藏层。注意，代价 $C_E(W,b,x^{(u)})$ 是激活量 $z_1^{(2)},\ldots,z_n^{(2)}$ 的函数，即
\begin{equation}
C_E(W,b,x^{(u)})=C_E(W,b,x^{(u)},z_1^{(2)},\ldots,z_n^{(2)}).
\tag{6.63}
\end{equation}
由链式法则，
\begin{equation}
\begin{aligned}
\delta_j^{(1)}(x^{(u)})&=\frac{\partial C_E(W,b,x^{(u)})}{\partial z_j^{(1)}}\\
&=\sum_{v=1}^n\frac{\partial C_E(W,b,x^{(u)},z_1^{(2)},\ldots,z_n^{(2)})}{\partial z_v^{(2)}}\frac{\partial z_v^{(2)}}{\partial z_j^{(1)}}\\
&=\sum_{v=1}^n\delta_v^{(2)}\frac{\partial z_v^{(2)}}{\partial z_j^{(1)}}.
\end{aligned}\tag{6.64}
\end{equation}
根据式 (6.27)、(6.30)，有
\begin{equation}
\frac{\partial z_v^{(2)}}{\partial z_j^{(1)}}=w_{vj}^{(2)}\sigma'(z_j^{(1)}).
\tag{6.65}
\end{equation}
将式 (6.65) 代入式 (6.64)，得到
\begin{equation}
\delta_j^{(1)}(x^{(u)})=\sum_{v=1}^n\delta_v^{(2)}w_{vj}^{(2)}\sigma'(z_j^{(1)}).
\tag{6.66}
\end{equation}
现在计算代价关于连接隐藏层的任意权重的变化率 $\partial C_E(W,b,x^{(u)})/\partial w_{jk}^{(1)}$。再次利用链式法则以及式 (6.26)、(6.64)，得到
\begin{equation}
\frac{\partial C_E(W,b,x^{(u)})}{\partial w_{jk}^{(1)}}=
\frac{\partial C_E(W,b,x^{(u)})}{\partial z_j^{(1)}}\frac{\partial z_j^{(1)}}{\partial w_{jk}^{(1)}}=\delta_j^{(1)}x_k^{(u)}.
\tag{6.67}
\end{equation}
下面计算代价关于网络中任意偏置的变化率 $\partial C_E(W,b,x^{(u)})/\partial b_j^{(l)}$。由式 (6.61)、(6.64)，得到
\begin{equation}
\frac{\partial C_E(W,b,x^{(u)})}{\partial b_j^{(2)}}=
\frac{\partial C_E(W,b,x^{(u)})}{\partial z_j^{(2)}}\frac{\partial z_j^{(2)}}{\partial b_j^{(2)}}=\delta_j^{(2)},
\tag{6.68}
\end{equation}
以及
\begin{equation}
\frac{\partial C_E(W,b,x^{(u)})}{\partial b_j^{(1)}}=
\frac{\partial C_E(W,b,x^{(u)})}{\partial z_j^{(1)}}\frac{\partial z_j^{(1)}}{\partial b_j^{(1)}}=\delta_j^{(1)}.
\tag{6.69}
\end{equation}
接着计算 $\partial C_{\mathrm{KL}}(p\Vert\widehat p)/\partial w_{jk}^{(l)}$ 与 $\partial C_{\mathrm{KL}}(p\Vert\widehat p)/\partial b_j^{(l)}$。利用式 (6.37)，得到
\begin{equation}
\begin{aligned}
\frac{\partial C_{\mathrm{KL}}(p\Vert\widehat p)}{\partial w_{jk}^{(l)}}
&=-\left[\frac p{\widehat p_j}-\frac{1-p}{1-\widehat p_j}\right]\frac{\partial\widehat p_j}{\partial w_{jk}^{(l)}}\\
&=\frac{\widehat p_j-p}{\widehat p_j(1-\widehat p_j)}\frac{\partial\widehat p_j}{\partial w_{jk}^{(l)}}.
\end{aligned}\tag{6.70}
\end{equation}
根据式 (6.36)，
\begin{equation}
\frac{\partial\widehat p_j(x^{(u)})}{\partial w_{jk}^{(l)}}=\sigma'(z_j^{(1)})x_k^{(u)}.
\tag{6.71}
\end{equation}

\bookdeepheading{paragraph}{6.2.1.3.2}{学习带非负约束的深度自编码器}

学习深度自编码器时，需要将单个自编码器的代价函数推广到多个自编码器。平均平方和误差 $C_{\mathrm{SAE}}(W,b,x)$ 不变，但需要修改用于度量稀疏性的 KL 距离。令
\begin{equation}
\widehat p_j^{(l)}=\frac1m\sum_{u=1}^ma_j^{(l)}(x^{(u)}),\quad j=1,\ldots,n_l,\ l=1,\ldots,L,
\tag{6.72}
\end{equation}
\begin{equation}
a_j^{(l)}(x^{(u)})=\sigma\left(\sum_{k=1}^{n_l}w_{jk}^{(l)}x_k^{(u)}+b_j^{(l)}\right).
\tag{6.73}
\end{equation}
对整个多重自编码器定义 KL 距离为
\begin{equation}
C_{\mathrm{SKL}}(P\Vert\widehat P)=\sum_{l=1}^L\sum_{j=1}^{n_l}
\left[p^{(l)}\log\frac{p^{(l)}}{\widehat p_j}+(1-p^{(l)})\log\frac{1-p^{(l)}}{1-\widehat p_j^{(l)}}\right].
\tag{6.74}
\end{equation}
首先计算 $\partial\widehat p_j^{(v)}/\partial w_{kg}^{(l)}$。对固定样本 $x^{(u)}$，考虑 $l>v$、$l=v$ 与 $l<v$ 三种情形。

（1）$l>v$。此时，$\widehat p_j^{(v)}$ 不是 $w_{kg}^{(l)}$ 的函数，因此
\begin{equation}
\frac{\partial\widehat p_j^{(v)}}{\partial w_{kg}^{(l)}}=0.
\tag{6.75}
\end{equation}
（2）$l=v$。根据式 (6.40)，得到
\begin{equation}
\frac{\partial\widehat p_j^{(v)}}{\partial w_{jg}^{(v)}}=\sigma'(z_j^{(v)})\frac{\partial z_j^{(v)}}{\partial w_{jg}^{(v)}}=
\sigma'(z_j^{(v)})a_g^{(v-1)},\quad g=1,\ldots,n_{v-1}.
\tag{6.76}
\end{equation}
（3）$l<v$。再次利用式 (6.40)，有
\begin{equation}
\begin{aligned}
\frac{\partial\widehat p_j^{(v)}}{\partial w_{qg}^{(l)}}&=\sigma'(z_j^{(v)})\sum_{k=1}^{n_{v-1}}w_{jk}^{(v)}\frac{\partial\widehat p_k^{(v-1)}}{\partial w_{qg}^{(l)}},\\
&\qquad q=1,\ldots,n_{v-1},\ q\ne j,\quad g=1,\ldots,n_{v-2},
\end{aligned}\tag{6.77}
\end{equation}
以及
\begin{equation}
\frac{\partial\widehat p_j^{(v)}}{\partial w_{jg}^{(l)}}=\sigma'(z_j^{(v)})w_{jg}^{(v)}\frac{\partial\widehat p_g^{(v-1)}}{\partial w_{jg}^{(l)}}.
\tag{6.78}
\end{equation}
类似地，关于 $\partial\widehat p_j^{(v)}/\partial b_j^{(l)}$，有
\begin{equation}
\frac{\partial\widehat p_j^{(v)}}{\partial b_g^{(l)}}=
\begin{cases}
0,&l>v,\\
\sigma'(z_j^{(v)}),&l=v,\\
\sigma'(z_j^{(v)})\displaystyle\sum_{k=1}^{n_{v-1}}w_{jk}^{(v)}
\frac{\partial\widehat p_k^{(v-1)}}{\partial b_g^{(l)}},\quad g=1,\ldots,n_{v-1},&l<v.
\end{cases}
\tag{6.79}
\end{equation}
综上，得到
\begin{equation}
\frac{\partial C_{\mathrm{SKL}}(p\Vert\widehat p)}{\partial w_{jk}^{(l)}}=
\sum_{v=l}^L\frac{\widehat p_j^{(v)}-p^{(v)}}{\widehat p_j^{(v)}(1-\widehat p_j^{(v)})}\frac{\partial\widehat p_j^{(v)}}{\partial w_{jk}^{(l)}},
\tag{6.80}
\end{equation}
\begin{equation}
\frac{\partial C_{\mathrm{SKL}}(p\Vert\widehat p)}{\partial b_g^{(l)}}=
\sum_{v=l}^L\frac{\widehat p_j^{(v)}-p^{(v)}}{\widehat p_j^{(v)}(1-\widehat p_j^{(v)})}\frac{\partial\widehat p_j^{(v)}}{\partial b_g^{(l)}}.
\tag{6.81}
\end{equation}
最后，利用式 (6.75)，得到
\begin{equation}
\frac{\partial C_{\mathrm{SW}}}{\partial w_{jk}^{(l)}}=g(w_{jk}^{(l)}),\quad j=1,\ldots,n_l,\ k=1,\ldots,n_{l-1},
\tag{6.82}
\end{equation}
其中
\[
g(w_{jk}^{(l)})=\begin{cases}|w_{jk}^{(l)}|,&w_{jk}^{(l)}<0,\\0,&w_{jk}^{(l)}\geq0.\end{cases}
\]
将上述结果合在一起，得到所需的代价函数偏导数：
\begin{equation}
\begin{aligned}
\frac{\partial C_{\mathrm{SAE}}(W,b)}{\partial w_{jk}^{(l)}}
&=\frac{\partial C_E(W,b)}{\partial w_{jk}^{(l)}}+\mu\frac{\partial C_{\mathrm{SKL}}(p\Vert\widehat p)}{\partial w_{jk}^{(l)}}+\lambda\frac{\partial C_{\mathrm{SW}}}{\partial w_{jk}^{(l)}}\\
&=\frac1{2m}\sum_{u=1}^m\delta_j^{(l)}(x^{(u)})a_k^{(l-1)}(x^{(u)})+,\\
&\quad\frac\mu m\sum_{v=l}^L\frac{\widehat p_j^{(v)}(x^{(u)})-p^{(v)}}{\widehat p_j^{(v)}(x^{(u)})(1-\widehat p_j^{(v)}(x^{(u)}))}
\frac{\partial\widehat p_j^{(v)}(x^{(u)})}{\partial w_{jk}^{(l)}}+\lambda g(w_{jk}^{(l)}),
\end{aligned}\tag{6.83}
\end{equation}
\begin{equation}
\begin{aligned}
\frac{\partial C_{\mathrm{SAE}}(W,b)}{\partial b_j^{(l)}}&=
\frac{\partial C_E(W,b)}{\partial b_j^{(l)}}+\mu\frac{\partial C_{\mathrm{SKL}}(p\Vert\widehat p)}{\partial b_j^{(l)}}\\
&=\frac1{2m}\sum_{u=1}^m\delta_j^{(l)}(x^{(u)})\\
&\quad+\frac\mu m\sum_{v=l}^L\frac{\widehat p_j^{(v)}(x^{(u)})-p^{(v)}}{\widehat p_j^{(v)}(x^{(u)})(1-\widehat p_j^{(v)}(x^{(u)}))}
\frac{\partial\widehat p_j^{(v)}(x^{(u)})}{\partial b_g^{(l)}}.
\end{aligned}\tag{6.84}
\end{equation}
学习深度自编码器的反向传播算法归纳为结论 6.2。

\Needspace{5\baselineskip}
\noindent\textbf{结论 6.2\quad 学习深度自编码器的反向传播算法}

\noindent\textbf{步骤 1：初始化。}预先指定较小的 $\varepsilon$，利用 $N(0,\varepsilon^2)$ 分布，将各层参数 $w_{jk}^{(l)}$ 与 $b_j^{(l)}$ 随机初始化为接近零的小值。

\noindent\textbf{步骤 2：进行前向激活计算。}

（1）计算输入层神经元的激活：
\[
z_j^{(1)}=\sum_{k=1}^nw_{jk}^{(1)}x_k+b_j^{(1)},\quad j=1,\ldots,n_1,
\qquad a_j^{(1)}=\sigma(z_j^{(1)}).
\]
（2）对于 $l=2,\ldots,L$，对第 $l$ 层的每个神经元，令
\[
z_j^{(l)}=\sum_{k=1}^{n_{l-1}}w_{jk}^{(l)}a_k^{(l-1)}+b_j^{(l)},\quad j=1,\ldots,n_l,
\qquad a_j^{(l)}=\sigma(z_j^{(l)}).
\]
\noindent\textbf{步骤 3：更新参数 $w_{jk}^{(l)}$ 与 $b_j^{(l)}$。}

（1）设定惩罚参数 $\mu$ 与 $\lambda$。

（2）对输出层的每个输出神经元 $j=1,\ldots,n$，以及第 $L-1$ 层的每个神经元 $k=1,\ldots,n_{L-1}$，令
\[
\begin{aligned}
\delta_j^{(L)}(x^{(u)})&=\frac{\partial C_E(W,b,x^{(u)})}{\partial z_j^{(L)}}\\
&=(x_j^{(u)}-a_j^L(x^{(u)}))\sigma'\bigl(z_j^{(L)}(x^{(u)}),\\
\frac{\partial C_E(W,b,x^{(u)})}{\partial w_{jk}^{(L)}}&=\delta_j^{(L)}(x^{(u)})a_k^{(L-1)},\\
\frac{\partial C_E(W,b,x^{(u)})}{\partial b_j^{(L)}}&=\delta_j^{(L)}(x^{(u)}),
\end{aligned}
\]
\[
g(w_{jk}^{(L)})=\begin{cases}|w_{jk}^{(L)}|,&w_{jk}^{(L)}<0,\\0,&w_{jk}^{(L)}\geq0,\end{cases}
\]
\[
\begin{aligned}
\frac{\partial C_{\mathrm{SAE}}(W,b)}{\partial w_{jk}^{(L)}}&=\frac1{2m}\sum_{u=1}^m\delta_j^{(L)}(x^{(u)})a_k^{(L-1)}(x^{(u)})+\lambda g(w_{jk}^{(L)}),\\
\frac{\partial C_{\mathrm{SAE}}(W,b)}{\partial b_j^{(L)}}&=\frac1{2m}\sum_{u=1}^m\delta_j^{(l)}(x^{(u)}),\\
w_{jk}^{(L)}&=w_{jk}^{(L)}-\alpha\frac{\partial C_{\mathrm{SAE}}(W,b)}{\partial w_{jk}^{(L)}},\\
b_j^{(L)}&=b_j^{(L)}-\alpha\frac{\partial C_{\mathrm{SAE}}(W,b)}{\partial b_j^{(L)}}.
\end{aligned}
\]
（3）对于 $l=L-1,\ldots,1$，
\[
\begin{aligned}
\delta_j^{(l)}(x^{(u)})&=\sum_{k=1}^{n_{l+1}}\delta_k^{(l+1)}(x^{(u)})w_{kj}^{(l+1)}\sigma'(z_j^{(l)}),\\
\frac{\partial C_E(W,b,x^{(u)})}{\partial w_{jk}^{(l)}}&=\delta_j^{(l)}(x^{(u)})a_k^{(l-1)}(x^{(u)}),\\
\frac{\partial C_E(W,b,x^{(u)})}{\partial b_j^{(l)}}&=\delta_j^{(l)}(x^{(u)}),\\
\widehat p_j^{(l)}&=\frac1m\sum_{u=1}^ma_j^{(l)}(x^{(u)}),\quad j=1,\ldots,n_l,\ l=1,\ldots,L,\\
a_j^{(l)}(x^{(u)})&=\sigma\left(\sum_{k=1}^{n_l}w_{jk}^{(l)}x_k^{(u)}+b_j^{(l)}\right),\\
\frac{\partial\widehat p_j^{(v)}}{\partial w_{kg}^{(l)}}&=0,\quad l>v,\\
\frac{\partial\widehat p_j^{(v)}}{\partial w_{jg}^{(v)}}&=\sigma'(z_j^{(v)})a_g^{(v-1)},\quad g=1,\ldots,n_{v-1},\ l=v,\\
\frac{\partial\widehat p_j^{(v)}}{\partial w_{qg}^{(l)}}&=\sigma'(z_j^{(v)})\sum_{k=1}^{n_{v-1}}w_{jk}^{(v)}\frac{\partial\widehat p_k^{(v-1)}}{\partial w_{qg}^{(l)}},\\
&\qquad q=1,\ldots,n_{v-1},\ q\ne j,\quad g=1,\ldots,n_{v-2},\ l<v,\\
\frac{\partial\widehat p_j^{(v)}}{\partial w_{jg}^{(l)}}&=\sigma'(z_j^{(v)})w_{jg}^{(v)}\frac{\partial\widehat p_g^{(v-1)}}{\partial w_{jg}^{(l)}},\quad l<v,
\end{aligned}
\]
\[
\frac{\partial\widehat p_j^{(v)}}{\partial b_g^{(l)}}=\begin{cases}
0,&l>v,\\
\sigma'(z_j^{(v)}),&l=v,\\
\sigma'(z_j^{(v)})\displaystyle\sum_{k=1}^{n_{v-1}}w_{jk}^{(v)}\frac{\partial\widehat p_k^{(v-1)}}{\partial b_g^{(l)}},\quad g=1,\ldots,n_{v-1},&l<v,
\end{cases}
\]
\[
\begin{aligned}
\frac{\partial C_{\mathrm{SKL}}(p\Vert\widehat p)}{\partial w_{jk}^{(l)}}&=\sum_{v=l}^L\frac{\widehat p_j^{(v)}-p^{(v)}}{\widehat p_j^{(v)}(1-\widehat p_j^{(v)})}\frac{\partial\widehat p_j^{(v)}}{\partial w_{jk}^{(l)}},\\
\frac{\partial C_{\mathrm{SKL}}(p\Vert\widehat p)}{\partial b_g^{(l)}}&=\sum_{v=l}^L\frac{\widehat p_j^{(v)}-p^{(v)}}{\widehat p_j^{(v)}(1-\widehat p_j^{(v)})}\frac{\partial\widehat p_j^{(v)}}{\partial b_g^{(l)}},\\
\frac{\partial C_{\mathrm{SW}}}{\partial w_{jk}^{(l)}}&=g(w_{jk}^{(l)}),\quad j=1,\ldots,n_l,\ k=1,\ldots,n_{l-1},
\end{aligned}
\]
其中
\[
g(w_{jk}^{(l)})=\begin{cases}|w_{jk}^{(l)}|,&w_{jk}^{(l)}<0,\\0,&w_j^{(l)}\geq0,\end{cases}
\]
\[
\begin{aligned}
\frac{\partial C_E(W,b)}{\partial w_{jk}^{(l)}}&=\frac1{2m}\sum_{u=1}^m\delta_j^{(l)}(x^{(u)})a_k^{(l-1)}(x^{(u)}),\\
\frac{\partial C_{\mathrm{SKL}}(p\Vert\widehat p)}{\partial w_{jk}^{(l)}}&=\frac1m\sum_{v=l}^L\frac{\widehat p_j^{(v)}(x^{(u)})-p^{(v)}}{\widehat p_j^{(v)}(x^{(u)})(1-\widehat p_j^{(v)}(x^{(u)}))}\frac{\partial\widehat p_j^{(v)}(x^{(u)})}{\partial w_{jk}^{(l)}},\\
\frac{\partial C_{\mathrm{SAE}}(W,b)}{\partial w_{jk}^{(l)}}&=\frac{\partial C_E(W,b)}{\partial w_{jk}^{(l)}}+\mu\frac{\partial C_{\mathrm{SKL}}(p\Vert\widehat p)}{\partial w_{jk}^{(l)}}+\lambda\frac{\partial C_{\mathrm{SW}}}{\partial w_{jk}^{(l)}},\\
\frac{\partial C_E(W,b)}{\partial b_j^{(l)}}&=\frac1{2m}\sum_{u=1}^m\delta_j^{(l)}(x^{(u)}),\\
\frac{\partial C_{\mathrm{SKL}}(p\Vert\widehat p)}{\partial b_j^{(l)}}&=\frac1m\sum_{v=l}^L\frac{\widehat p_j^{(v)}(x^{(u)})-p^{(v)}}{\widehat p_j^{(v)}(x^{(u)})(1-\widehat p_j^{(v)}(x^{(u)}))}\frac{\partial\widehat p_j^{(v)}(x^{(u)})}{\partial b_g^{(l)}},\\
\frac{\partial C_{\mathrm{SAE}}(W,b)}{\partial b_j^{(l)}}&=\frac{\partial C_E(W,b)}{\partial b_j^{(l)}}+\mu\frac{\partial C_{\mathrm{SKL}}(p\Vert\widehat p)}{\partial b_j^{(l)}},\\
w_{jk}^{(l)}&=w_{jk}^{(l)}-\alpha\frac{\partial C_{\mathrm{SAE}}(W,b)}{\partial w_{jk}^{(l)}},\\
b_j^{(l)}&=b_j^{(l)}-\alpha\frac{\partial C_{\mathrm{SAE}}(W,b)}{\partial b_j^{(l)}}.
\end{aligned}
\]
\noindent\textbf{步骤 4：检查收敛性。}若 $|C_{\mathrm{SAE}}(W^{(t+1)},b^{(t+1)})-|C_{\mathrm{SAE}}(W^{(t)},b^{(t)})|<\varepsilon$，即当前迭代与前一次迭代之间的代价函数差小于预先指定的误差 $\varepsilon$，则停止；否则，返回步骤 1，重复迭代。
